\section{总结与展望}\label{sec:summary}

\subsection{本文工作总结}

本文围绕最小哈希族的显式构造问题，系统研究了在种子长度与乘法误差之间的权衡关系，并从有限次独立性方法、结构化伪随机构造以及随机性提取器三个层面给出了较为完整的理论刻画。

首先，\ref{sec:bounded_independence} 在有限次独立性框架下，我们精确刻画了实现 $k$-最小哈希所需的独立次数。通过上界分析与下界构造的匹配，证明了 $\Theta(k + \log 1/\delta)$ 次独立性是充分且必要的，从而建立了独立性与误差之间的最优权衡关系。同时，对经典仿射哈希族的分析表明，代数结构所带来的线性相关性会导致不可避免的乘法误差下界，这说明仅依赖有限次独立性难以在小误差情形下达到最优种子长度。

在此基础上，本文进一步研究了如何突破有限次独立性方法的局限。\ref{sec:optimal_seed_length} 的核心贡献是在最优种子长度 $O(\log N)$ 下，实现了亚常数级别的乘法误差，并将该结果推广至 $k$-最小哈希情形。该构造的关键在于避免将加法误差直接转化为乘法误差所带来的损失，而是通过结构化分解对概率进行精细控制。具体而言，我们利用分桶哈希将全局事件分解为多个局部事件的乘积结构，通过单次只读分支程序刻画桶间关系，并结合单层 Nisan-Zuckerman 伪随机生成器实现随机性的跨桶复用，从而在压缩种子长度的同时保持整体分布的可控性。

在误差分析中，一个关键困难来源于小概率事件对乘法误差的放大效应。为此，我们引入具有均匀性性质的提取器，使得关键部分的分布能够被精确刻画，而其余部分则通过组合矩形伪随机生成器及其域缩减技术进行控制。进一步地，在推广到 $k$-最小哈希情形时，通过与有限次独立哈希族的直和构造，对涉及目标元素的小概率事件进行单独处理，从而在整体结构上实现误差与种子长度的兼顾。

最后，\ref{sec:extractor} 从随机性提取器的角度对上述构造中的关键技术进行了系统分析。一方面，我们构造了具有均匀性性质的提取器，并证明其可以通过线性结构与正交分解在递归过程中保持该性质；另一方面，我们引入乘法误差提取器的概念并证明其在一般意义下不可实现。这一结果表明，最小哈希问题中的乘法误差控制无法通过简单强化提取器模型获得，而必须依赖于结构性的概率分解与精细分析。

总体而言，本文的主要贡献在于：在最优种子长度框架下，建立了一套处理乘法误差的系统方法，将分桶结构、伪随机生成器与提取器技术有机结合，从而在理论上填补了最小哈希族构造中的关键空白。

\subsection{讨论与展望}

我们首先讨论一些在正文中省略的技术细节：
\begin{enumerate}
    \item 第一个问题是有关直和结构的适用性。设 $G_1(s_1)$ 是函数族 $\mathcal{F}_1 = \{f_1 : [K] \to \{0, 1\}\}$ 误差为 $\varepsilon_1$ 的伪随机生成器，$G_2(s_2)$ 是函数族 $\mathcal{F}_2 = \{f_2 : [K] \to \{0, 1\}\}$ 误差为 $\varepsilon_2$ 的伪随机生成器。考虑其直和
    $$
    G(s_1, s_2) = (G_1 \oplus G_2)(s_1, s_2) := G_1(s_1) + G_2(s_2),
    $$
    一个自然的问题是 $G$ 是否继承 $G_1$ 与 $G_2$ 的伪随机性性质。

    该结论并非普适成立，其成立的关键条件是函数族 $\mathcal{F}_1$ 与 $\mathcal{F}_2$ 需满足\emph{平移不变性}。具体而言，若对任意 $b \in [K]$ 与 $f \in \mathcal{F}$，定义 $\mathsf{shift}_b(f) : x \mapsto f(x + b)$，并且始终有 $\mathsf{shift}_b(f) \in \mathcal{F}$，则称 $\mathcal{F}$ 为平移不变的。

    在此条件下，对于任意 $f_1 \in \mathcal{F}_1$ 与任意 $b \in [K]$，由 $G_1$ 的伪随机性有
    $$
    \E_{s_1}[\mathsf{shift}_b(f_1)(G_1(s_1))] 
    = \E[\mathsf{shift}_b(f_1)(U)] \pm \varepsilon_1
    = \E[f_1(U + b)] \pm \varepsilon_1
    = \E[f_1(U)] \pm \varepsilon_1.
    $$
    进一步利用 $s_1$ 与 $s_2$ 的独立性，对任意 $f_1 \in \mathcal{F}_1$ 有
    \begin{align*}
        \E_{s_1, s_2}[f_1(G(s_1, s_2))] 
        &= \sum_{b \in [K]} \E_{s_1}[f_1(G_1(s_1) + b)] \cdot \Pr_{s_2}[G_2(s_2) = b] \\
        &= \sum_{b \in [K]} \E_{s_1}[\mathsf{shift}_b(f_1)(G_1(s_1))] \cdot \Pr_{s_2}[G_2(s_2) = b] \\
        &= \sum_{b \in [K]} \left( \E[f_1(U)] \pm \varepsilon_1 \right) \cdot \Pr_{s_2}[G_2(s_2) = b] \\
        &= \E[f_1(U)] \pm \varepsilon_1.
    \end{align*}
    同理可得 $G_2$ 的性质亦可保持，从而 $G$ 同时继承两个伪随机生成器的性质。

    特别地，$t$-集团函数类与组合矩形函数类均满足平移不变性。因此在 \ref{sec:optimal_seed_length} 中最小哈希族构造第 $5$ 步得到的直和族 $\mathcal{G}$ 同时满足所需的两个伪随机性性质。

    \item 第二个问题是关于为何 \ref{sec:optimal_seed_length} 中采用的是单层 Nisan-Zuckerman 伪随机生成器，而非 Nisan 伪随机生成器或 Impagliazzo-Nisan-Wigderson 伪随机生成器。

    其主要原因在于后两类构造均基于递归结构，因此对输入顺序具有较强依赖性，无法进行任意重排 \cite{tzur2009notions}。这一限制使得在处理包含 $y$ 的特殊桶时较为困难，而该部分往往对应最大的熵损失。

    此外，目前尚不清楚如何将组合矩形伪随机生成器中的域缩减技术有效地应用到 Nisan 或 Impagliazzo-Nisan-Wigderson 伪随机生成器中。因此，在本工作的构造中采用单层 Nisan-Zuckerman 伪随机生成器更为合适。
\end{enumerate}

最后讨论若干开放问题与未来研究方向：
\begin{enumerate}
    \item \textbf{第一个问题是，如何在保持种子长度的同时进一步减小乘法误差。}
    
    对于最小哈希族，一个重要目标是在实现 $O(\log N)$ 种子长度的同时，将误差降低至多项式级别。一个可能的方向是将常数次独立哈希族与组合矩形伪随机生成器进行直和，并分析
    $$
    \Pr_h[h(y) < \min h(X)] 
    = \sum_\theta \Pr_h[h(y) = \theta] \cdot \Pr_h[\min h(X) > \theta \mid h(y) = \theta].
    $$

    其中 $\Pr_h[h(y) = \theta]$ 可由常数次独立性较好地控制。然而，条件概率项目前尚无法类似 \cref{lem:PRG_CR_LOG} 的方式进行有效分析。由于 $\Pr_{\sigma \sim U}[\sigma(y) = \theta] = 1/M$，为了保证
    $\Pr_h[h(y) = \theta] > 0$ 从而避免条件概率退化，必须要求伪随机生成器的误差满足 $\varepsilon < 1/M$，这将导致种子长度回到 $O(\log N \log \log N)$ 的量级。

    因此，该直和构造是否存在更巧妙的分析方法，或其本身是否存在内在限制，仍有待进一步研究。

    对于 $k$-最小哈希族，基于 \ref{sec:bounded_independence} 有限次独立性的构造已经可以在 $k = \Omega(\log N)$ 时实现 $1/\poly(N)$ 的误差。然而在 $k = o(\log N)$ 的情形下，如何达到多项式级误差仍不清楚。此外，由于 $k$-最小哈希的最优种子长度为 $O(k \log N + \log 1 / \delta)$，当 $\delta$ 极小（例如 $\delta \le 1/N^k$）时，是否仍可达到 $O(k \log N)$ 的种子长度也是一个重要的开放问题。

    \item \textbf{第二个问题是，如何降低哈希函数的计算时间（evaluation time）。}
    
    本文中的构造均为显式的，即具有多项式时间的计算复杂度，但未对具体计算效率作进一步优化。

    基于有限次独立性的构造通常需要 $O(\log 1/\delta)$ 次字操作，而组合矩形伪随机生成器可实现 $O((\log \log N)^2)$ 次字操作。相比之下，\ref{sec:optimal_seed_length} 中的分桶哈希构造在计算效率方面仍存在改进空间。

    因此，在保持理论性质的同时进一步优化计算复杂度，是一个具有实际意义的重要研究方向。
\end{enumerate}